%! TeX program = lualatex

\documentclass[12pt]{article}

\usepackage{cmap}

\usepackage[english, russian]{babel}

\usepackage{float}
\usepackage{fontspec}
\usepackage{libertinus}
\setmonofont{Source Code Pro}

\usepackage{microtype}
\usepackage{graphicx}
\usepackage{xcolor} % Цвета

\usepackage{subfig}

\usepackage{listings}

\usepackage[left=2.5cm, right=2.5cm, top=2cm, bottom=2cm]{geometry}

\usepackage[mocha, textcolor=true, pagecolor=true]{catppuccinpalette}
% \usepackage[latte, textcolor=true, pagecolor=true]{catppuccinpalette}


\definecolor{codegreen}{rgb}{0,0.6,0}
\definecolor{codegray}{rgb}{0.5,0.5,0.5}
\definecolor{codepurple}{rgb}{0.796, 0.651, 0.969}
\definecolor{codeBlue}{rgb}{0.537, 0.706, 0.980}

\lstdefinestyle{mystyle}{
    commentstyle=\color{codegreen},
    keywordstyle=\color{codeBlue},
    numberstyle=\tiny\color{codegray},
    stringstyle=\color{codepurple},
    basicstyle=\ttfamily\footnotesize,
    breakatwhitespace=false,         
    frame=leftline,
    xleftmargin=8pt,
    breaklines=true,                 
    captionpos=b,                    
    keepspaces=true,                 
    numbers=left,                    
    numbersep=5pt,                  
    showspaces=false,                
    showstringspaces=false,
    showtabs=false,                  
    tabsize=2
}

\lstset{style=mystyle}
\renewcommand{\lstlistingname}{}

\begin{document}

\begin{center}
	\section*{Лабораторная работа №3}
	Вариант №3
\end{center}

\noindent Листинг библиотеки классов (dll):

\begin{lstlisting}[language=Python]
using System;

namespace dll
{
    public class Machine
    {
        public string Metal { get; set; } = "None";
        public int Age { get; set; } = 0;
        public bool Break { get; set; } = false;

        public void Repair()
        {
            if (Break)
            {
                Break = false;
                Console.WriteLine("Car is fixed");
            }

            else
                Console.WriteLine("Car is in good condition");

        }

        public void setAge(int new_age)
        {
            Age = new_age;
            Console.WriteLine($"Set age: {Age}");
        }

        public void ShowInfo()
        {
            string Status = Checking();
            Console.WriteLine($"\nMachine age: {Age}\nStatus: {Status}\nMetal: {Metal}");
        }

        public string Checking()
        {
            if (!Break)
                return "serviceable";
            else
                return "break";
        }
    }

    public class Car : Machine
    {
        public int id { get; set; } = 0;
        public string Model { get; set; } = "None";
        public int Places = 0;
        public Engine engine = new Engine();
        public bool Drive = false;


        public void DriveToggle()
        {
            if (Drive) { Drive = false; }
            else { Drive = true; }

            if (Drive && !Break)
            {


                if (!engine.Work)
                {
                    engine.Start();
                }

                Console.WriteLine("Drive");
            }

            else if (Break)
            {
                Drive = false;
                Console.WriteLine("Car is broken!");
            }
            else
            {
                if (engine.Work)
                {
                    engine.Stop();
                }
                Console.WriteLine("Stop drive");
            }
        }

        public void Honk() { Console.WriteLine("HONK!"); }

        public void ReplaceEngine()
        {
            if (engine.Work) { engine.Stop(); }

            Console.WriteLine("Model engine: ");
            string model = Console.ReadLine();

            Console.WriteLine("Cylinders engine: ");
            int cylinders = int.Parse(Console.ReadLine());

            Console.WriteLine("HorsePower engine: ");
            int horse_power = int.Parse(Console.ReadLine());

            Engine new_engine = new Engine { model_engine = model, Cylinders = cylinders, HorsePower = horse_power };
            engine = new_engine;

        }

        public void modelEngine()
        {
            Console.WriteLine(engine.model_engine);

        }

        public Car CreateCar()
        {
            Car newCar = new Car();

            Console.WriteLine("Model: ");
            newCar.Model = Console.ReadLine();

            Console.WriteLine("Places: ");
            newCar.Places = int.Parse(Console.ReadLine());

            Console.WriteLine("--Machine Properties--");
            Console.WriteLine("Metal: ");
            newCar.Metal = Console.ReadLine();

            Console.WriteLine("Age: ");
            newCar.Age = int.Parse(Console.ReadLine());

            Console.WriteLine("Break (true/false): ");
            newCar.Break = bool.Parse(Console.ReadLine());

            Console.WriteLine("--Engine Properties--");
            /*newCar.engine = new Engine();*/
            newCar.engine.createEngine();

            return newCar;

        }

    }

    public class Engine
    {
        public int id { get; set; } = 0;
        public string model_engine { get; set; } = "None";
        public int Cylinders { get; set; } = 0;
        public int HorsePower { get; set; } = 0;
        public bool Work = false;

        public Engine createEngine()
        {

            Console.WriteLine("Model: ");
            string Model_engine = Console.ReadLine();
            Console.WriteLine("Cylinders: ");
            int cylinders = int.Parse(Console.ReadLine());
            Console.WriteLine("HorsePower: ");
            int horsePower = int.Parse(Console.ReadLine());

            return new Engine { Cylinders = Cylinders, HorsePower = horsePower, model_engine = Model_engine };
        }

        public void Start()
        {
            if (!Work) { Work = true; } else { Console.WriteLine("The engine is already running"); }
        }

        public void Stop()
        {
            if (Work) { Work = false; } else { Console.WriteLine("The engine is no longer running"); }
        }


        public void Upgrade(int extraHorsPower)
        {
            HorsePower = extraHorsPower;
        }

    }
}
\end{lstlisting}

\vspace{12pt}

\noindent Листинг веб-сервиса (asmx.cs):

\begin{lstlisting}[language=Python]
using dll;
using System;
using System.ComponentModel;
using System.Configuration;
using System.Data;
using System.Data.SqlClient;
using System.Web.Services;

namespace freakWeb
{

[WebService(Namespace = "http://tempuri.org/")]
	[WebServiceBinding(ConformsTo = WsiProfiles.BasicProfile1_1)]
	[System.ComponentModel.ToolboxItem(false)]


public class WebService1 : System.Web.Services.WebService
{

private string connectionString = ConfigurationManager.ConnectionStrings["db"].ConnectionString;

[WebMethod]
public int CreateCars(string metal, int age, string model, bool @break, int places)
{
		using (SqlConnection connection = new SqlConnection(connectionString))
		{
				connection.Open();
				string sqlInsertMachine = "INSERT INTO Machines (Metal, Age, [Break]) VALUES (@Metal, @Age, @Break); SELECT SCOPE_IDENTITY();";
				int machineId;
				using (SqlCommand command = new SqlCommand(sqlInsertMachine, connection))
				{
						command.Parameters.Add("@Metal", SqlDbType.NVarChar, 50).Value = metal;
						command.Parameters.Add("@Age", SqlDbType.Int).Value = age;
						command.Parameters.Add("@Break", SqlDbType.Bit).Value = @break;
						machineId = Convert.ToInt32(command.ExecuteScalar());
					}


				string sql = "INSERT INTO Cars (Model, Places, Id) VALUES (@Model, @Places, @Id); SELECT SCOPE_IDENTITY();";

				using (SqlCommand command = new SqlCommand(sql, connection))
				{
						command.Parameters.Add("@Id", SqlDbType.Int).Value = machineId;
						command.Parameters.Add("@Model", SqlDbType.NVarChar, 50).Value = model;
						command.Parameters.Add("@Places", SqlDbType.Int).Value = places;
						command.ExecuteNonQuery();
						return machineId;
					}
			}
	}

	[WebMethod]
public BindingList<Car> GetAllCars()
{
		BindingList<Car> cars = new BindingList<Car>();

		using (SqlConnection connection = new SqlConnection(connectionString))
		{
				connection.Open();
				string sql = "SELECT * FROM Cars";

				using (SqlCommand command = new SqlCommand(sql, connection))
				{
						using (SqlDataReader reader = command.ExecuteReader())
						{
								while (reader.Read())
								{
										Car car = new Car
											{
												id = (int)reader["Id"],
												Model = reader["Model"] == DBNull.Value ? null : (string)reader["Model"],
												Places = reader["Places"] == DBNull.Value ? 0 : (int)reader["Places"],
												Drive = reader["Drive"] == DBNull.Value ? false : (bool)reader["Drive"],

											};
										cars.Add(car);
									}
							}
					}
			}
		return cars;
	}
}
}
\end{lstlisting}

\vspace{12pt}

\noindent Листинг приложения Windows Forms

\begin{lstlisting}[language=Python]
using Forms.ServiceReference1;
using System;
using System.ComponentModel;
using System.Drawing;
using System.IO;
using System.Windows.Forms;
using System.Xml.Serialization;

namespace Forms
{
    public partial class Form1 : Form
    {
        private BindingList<Car> car_list = new BindingList<Car>();
        public int count_car = 1;
        public Form1()
        {
            InitializeComponent();
            // dataGridView1.DataSource = car_list;
            comboBox1.Items.AddRange(new object[] { "id", "Metal", "Age", "Break" });
        }

        private void Search_Click_1(object sender, EventArgs e)
        {
            string searchValue = textBox1.Text;
            string chouse = comboBox1.Text;

            foreach (DataGridViewRow row in dataGridView1.Rows)
            {
                row.DefaultCellStyle.BackColor = Color.White; // сброс цвета

                if (row.Cells[chouse].Value?.ToString() == searchValue)
                {
                    row.DefaultCellStyle.BackColor = Color.Tan;
                }
            }
        }

        private void ClearShearch_Click(object sender, EventArgs e)
        {
            foreach (DataGridViewRow row in dataGridView1.Rows)
            {
                row.DefaultCellStyle.BackColor = Color.White;
            }
        }

        private void CreateCustomCar_Click_1(object sender, EventArgs e)
        {
            bool Break = inputBreak.Checked;

            string Model = inputModel.Text;
            string Metal = inputMetal.Text;
            int Age = int.Parse(inputAge.Text);
            int places = int.Parse(inputPlaces.Text);

            using (WebService1SoapClient client = new WebService1SoapClient())
            {
                client.CreateCars(Metal, Age, Model, Break, places);
            }
        }
        private void SaveCarToXml(BindingList<Car> car, string filePath)
        {
            XmlSerializer serializer = new XmlSerializer(typeof(BindingList<Car>));
            using (StreamWriter writer = new StreamWriter(filePath))
            {
                serializer.Serialize(writer, car);
            }
        }

        private BindingList<Car> LoadCarFromXml(string filePath)
        {
            XmlSerializer serializer = new XmlSerializer(typeof(BindingList<Car>));
            using (StreamReader reader = new StreamReader(filePath))
            {
                return (BindingList<Car>)serializer.Deserialize(reader);
            }
        }

        private void Grid2xml_Click_1(object sender, EventArgs e)
        {
            SaveFileDialog saveDialog = new SaveFileDialog();
            if (saveDialog.ShowDialog() == DialogResult.OK)
            {
                SaveCarToXml(car_list, saveDialog.FileName);
            }
        }

        private void Xml2grid_Click_1(object sender, EventArgs e)
        {
            OpenFileDialog openDialog = new OpenFileDialog();
            if (openDialog.ShowDialog() == DialogResult.OK)
            {
                car_list = LoadCarFromXml(openDialog.FileName);
                dataGridView1.DataSource = car_list;
            }
        }

        private void CreateTestCar_Click_1(object sender, EventArgs e)
        {
            Car car = new Car
            {
                id = count_car,
                Model = "Tesla",
                Metal = "Aluminum",
                Age = 2,
                Break = false,
            };
            car_list.Add(car);
            count_car++;
        }


        private async void Sql2grid_Click(object sender, EventArgs e)
        {
            WebService1SoapClient client = new WebService1SoapClient();

            ServiceReference1.GetAllCarsResponse response = await client.GetAllCarsAsync();

            BindingList<ServiceReference1.Car> cars = new BindingList<ServiceReference1.Car>(response.Body.GetAllCarsResult);

            dataGridView1.DataSource = cars;

            if (client.State != System.ServiceModel.CommunicationState.Faulted)
            {
                client.Close();
            }
        }

    }
}
\end{lstlisting}

\vspace{12pt}

\noindent Листинг SQL скриптов:

\begin{itemize}
	\item cars.sql

	      \begin{lstlisting}[language = Python]
CREATE TABLE [dbo].[Car] (
    [Id] INT NOT NULL PRIMARY KEY,
    [Model] NVARCHAR(10) NULL DEFAULT 'None',
);


\end{lstlisting}

	\item engine.sql

	      \begin{lstlisting}[language = Python]
CREATE TABLE [dbo].[Engines] (
    [Id] INT NOT NULL PRIMARY KEY IDENTITY(1,1),
    [ModelEngine] NVARCHAR(10) NULL DEFAULT 'None', 
    [Cylinders] INT NULL DEFAULT 0,
    [HorsePower] INT NULL DEFAULT 0,
    [Work] BIT NULL DEFAULT 0
);

\end{lstlisting}

	      \break

	\item machine.sql

	      \begin{lstlisting}[language = Python ]
CREATE TABLE [dbo].[Machines]
(
    [Id] INT NOT NULL PRIMARY KEY IDENTITY(1,1),
    [Metal] NVARCHAR(10) NULL DEFAULT 'None',   
    [Age] INT NULL DEFAULT 0,
    [Break] BIT NULL DEFAULT 0 
)
\end{lstlisting}

\end{itemize}

\subsection*{Результаты работы программы:}

\end{document}
